\section*{第 \ref{ch:g10:inertia} 章——力与惯性原理}

\begin{solution}{exo:g10:inertia:1}
(a) \omterm{def:g10:universal-gravitation:weight}{重力}（\omterm{def:g10:inertia:contact}{非接触力}），桌面的\omterm{def:g10:inertia:inventory}{法向力}（\omterm{def:g10:inertia:contact}{接触力}）。
(b) \omterm{def:g10:universal-gravitation:weight}{重力}（\omterm{def:g10:inertia:contact}{非接触力}），细线的\omterm{def:g10:inertia:inventory}{张力}（\omterm{def:g10:inertia:contact}{接触力}）。
(c) \omterm{def:g10:universal-gravitation:weight}{重力}（\omterm{def:g10:inertia:contact}{非接触力}），冰面的\omterm{def:g10:inertia:inventory}{法向力}（\omterm{def:g10:inertia:contact}{接触力}）；\omterm{def:g10:inertia:inventory}{摩擦力}
可忽略。(d) 只有\omterm{def:g10:universal-gravitation:weight}{重力}——没有物体接触苹果。
\end{solution}

\begin{solution}{exo:g10:inertia:2}
书：$P = 0.250 \times 9.81 \approx \qty{2.45}{N}$，箭头长约
\qty{2.5}{cm}。人：$P = 60 \times 9.81 \approx \qty{589}{N}$
——箭头长 \qty{5.89}{m}。选取合适的比例尺，
使最大的力恰能画在纸上。
\end{solution}

\begin{solution}{exo:g10:inertia:3}
\omterm{def:g10:inertia:force}{作用点} $A$；方向水平，指向左侧；
大小 $3.5 \times 2 = \qty{7.0}{N}$。
\end{solution}

\begin{solution}{exo:g10:inertia:4}
\omterm{def:g10:universal-gravitation:weight}{重力} $\vect P$：竖直向下，$P = 1.2 \times 9.81 \approx
\qty{11.8}{N}$。架子的\omterm{def:g10:inertia:inventory}{法向力} $\vect R$：与\omterm{def:g10:universal-gravitation:weight}{重力}在同一直线上，
方向竖直向上，$R = \qty{11.8}{N}$（字典静止，因此二者
\omterm{def:g10:inertia:compensate}{相互抵消}）。
\end{solution}

\begin{solution}{exo:g10:inertia:5}
(a) 错误：\omterm{def:g10:inertia:inventory}{摩擦力}使物体停下来；若不受任何力，物体将保持
\omterm{def:g10:inertia:uniform}{匀速直线运动}。(b) 正确（惯性原理）。
(c) 错误：它所受的力\emph{\omterm{def:g10:inertia:compensate}{相互抵消}}（例如\omterm{def:g10:universal-gravitation:weight}{重力}
与\omterm{def:g10:inertia:inventory}{法向力}）。(d) 正确——这是惯性原理的逆否
表述。
\end{solution}

\begin{solution}{exo:g10:inertia:6}
\emph{1.} $v = 24/3.0 = \qty{8.0}{m/s}$。受力为\omterm{def:g10:universal-gravitation:weight}{重力}与\omterm{def:g10:inertia:inventory}{法向力}，
它们\omterm{def:g10:inertia:compensate}{相互抵消}（\omterm{def:g10:inertia:uniform}{匀速直线运动}，\omterm{def:g10:inertia:inventory}{摩擦力}
可忽略）。\emph{2.} 混凝土地面的\omterm{def:g10:inertia:inventory}{摩擦力}不再
可忽略：冰球受到一个未被抵消的向后的力，
因而速度确实改变了——冰球减速并停止。
\end{solution}

\begin{solution}{exo:g10:inertia:7}
\omterm{def:g10:universal-gravitation:weight}{重力} $P = 70 \times 9.81 \approx \qty{687}{N}$，向下；地面的
作用力向上。\omterm{def:g10:inertia:uniform}{匀速直线运动}，二者\omterm{def:g10:inertia:compensate}{相互抵消}：地面以
\qty{687}{N} 的力推人。静止时没有变化——静止与匀速运动
给出相同的\omterm{def:g10:inertia:diagram}{受力图}。
\end{solution}

\begin{solution}{exo:g10:inertia:8}
\emph{1.} 静止时，\omterm{def:g10:inertia:inventory}{张力}与\omterm{def:g10:universal-gravitation:weight}{重力}\omterm{def:g10:inertia:compensate}{相互抵消}：
$T = 1.8 \times 9.81 \approx \qty{17.7}{N}$。\emph{2.} 每根竖直
细线承受一半：$T = \qty{8.8}{N}$。
\end{solution}

\begin{solution}{exo:g10:inertia:9}
刹车时：乘客的身体因惯性保持原有速度，而
汽车则减速，因此乘客\emph{相对于
汽车}向前运动。转弯时：身体倾向于继续沿直线运动，而车身
在其下方转弯，因此乘客向弯道外侧漂移。两种
情况中所谓的“力”都是虚构的；真正缺少的是一个实际的力
（抓握力、扶手）来使乘客的速度随
汽车一起改变。
\end{solution}

\begin{solution}{exo:g10:inertia:10}
\omterm{def:g10:inertia:uniform}{匀速直线运动}，\omterm{def:g10:universal-gravitation:weight}{重力}与空气阻力\omterm{def:g10:inertia:compensate}{相互抵消}：阻力
$= P = 85 \times 9.81 \approx \qty{834}{N}$，竖直向上。
\end{solution}

\begin{solution}{exo:g10:inertia:11}
\emph{1.} \omterm{def:g10:inertia:uniform}{匀速直线运动}：水平方向的力\omterm{def:g10:inertia:compensate}{相互抵消}，因此
阻力为 \qty{300}{N}，方向与绳子相反。\emph{2.} 滑雪者的\omterm{def:g10:universal-gravitation:weight}{重力}
与水对滑雪板向上的\omterm{def:g10:inertia:inventory}{推力}。
\end{solution}

\begin{solution}{exo:g10:inertia:12}
\emph{1.} $P = 2.4 \times 9.81 \approx \qty{23.5}{N}$；竖直分力
承担\omterm{def:g10:universal-gravitation:weight}{重力}：$2T\cos 40^\circ = P$，因此
$T = 23.5/(2 \times 0.766) \approx \qty{15.4}{N}$。\emph{2.}
$T = P/(2\cos\theta)$ 随 $\theta \to 90^\circ$ 无限增大，
因为 $\cos\theta \to 0$；恰好水平的细线\emph{没有}
竖直分量，因此没有任何力能承担\omterm{def:g10:universal-gravitation:weight}{重力}。
\end{solution}

\begin{solution}{exo:g10:inertia:13}
推：\omterm{def:g10:universal-gravitation:weight}{重力}、\omterm{def:g10:inertia:inventory}{法向力}、手的\omterm{def:g10:inertia:inventory}{推力}、较小的\omterm{def:g10:inertia:inventory}{摩擦力}；速率
增大，因此这些力不\omterm{def:g10:inertia:compensate}{相互抵消}
（见\cref{prop:g10:inertia:converse}）。滑行：\omterm{def:g10:universal-gravitation:weight}{重力}、\omterm{def:g10:inertia:inventory}{法向力}、较小的\omterm{def:g10:inertia:inventory}{摩擦力}；
石块减速，因此同样不\omterm{def:g10:inertia:compensate}{相互抵消}——\omterm{def:g10:inertia:inventory}{摩擦力}是
未被抵消的剩余力（在理想冰面上将永远滑行）。静止：
仅受\omterm{def:g10:universal-gravitation:weight}{重力}与\omterm{def:g10:inertia:inventory}{法向力}，二者\omterm{def:g10:inertia:compensate}{相互抵消}
（见\cref{thm:g10:inertia:principle}）。
\end{solution}

\begin{solution}{exo:g10:inertia:14}
\emph{1.} 由惯性原理，保持恒定速度
无需任何力——而在星际空间中几乎没有力
作用。亚里士多德会预言，探测器在“推动者”
\emph{2.} 一年 $\approx \qty{3.156e7}{s}$，因此
$d = 17 \times \num{3.156e7} \approx \qty{5.4e8}{km} \approx
\qty{3.6}{au}$。\emph{3.} 速度是矢量：改变其方向
即改变速度，这需要一个未被抵消的力。
\end{solution}

\begin{solution}{exo:g10:inertia:15}
静止时三个力\omterm{def:g10:inertia:compensate}{相互抵消}：
$\vect F_3 = -(\vect F_1 + \vect F_2)$，因此 $\vect F_3\,(-2, -4)$，
大小 $\sqrt{4 + 16} = \sqrt{20} \approx \qty{4.5}{N}$，方向
与前两力的合力相反（网格中为左下方）。
\end{solution}

\begin{solution}{pb:g10:inertia:1}
\textbf{1.} \omterm{def:g10:inertia:uniform}{匀速直线运动}，速度为 \qty{0.75}{m/s}。受力：\omterm{def:g10:universal-gravitation:weight}{重力}
和传送带的\omterm{def:g10:inertia:inventory}{法向力}（无需\omterm{def:g10:inertia:inventory}{摩擦力}：行李箱
不在传送带上滑动）。
\textbf{2.} 是——\omterm{def:g10:inertia:uniform}{匀速直线运动}
（见\cref{thm:g10:inertia:principle}）。$P = 23 \times 9.81 \approx
\qty{226}{N}$，因此 $R = \qty{226}{N}$。
\textbf{3.} 两图相同。静止与\omterm{def:g10:inertia:uniform}{匀速直线运动}
是同一种力学状态：惯性
不能将二者区分开。
\textbf{4.} 传送带停止时，行李箱仍保持 \qty{0.75}{m/s} 的速度，
它向前滑动或倾倒，直至\omterm{def:g10:inertia:inventory}{摩擦力}使其停下。
没有任何东西推它——它只是保持了原有的速度。
\textbf{5.} 沿同一方向的速率相加：$0.75 + 1.2 =
\qty{1.95}{m/s}$。她的速度恒定，因此作用在她身上的
力\omterm{def:g10:inertia:compensate}{相互抵消}。
\textbf{6.} $P = 0.160 \times 9.81 \approx \qty{1.57}{N}$，方向向下，
$R = \qty{1.57}{N}$，方向向上，且\emph{无}水平方向的力。
\textbf{7.} 冰球以 \qty{8.0}{m/s} 的速度运动，身后没有任何推动者。
日常表面掩盖了\omterm{def:g10:inertia:inventory}{摩擦力}——一个未被抵消的反向力——
因此运动似乎会自行停止。
\textbf{8.} 其速度发生改变，因此作用力不再\omterm{def:g10:inertia:compensate}{相互抵消}
（见\cref{prop:g10:inertia:converse}）；原因是
粗糙冰面的\omterm{def:g10:inertia:inventory}{摩擦力}。
\textbf{9.} 不是：速度是\emph{矢量}，其方向
在改变，因此作用力不能\omterm{def:g10:inertia:compensate}{相互抵消}——冰面
对她的冰刀施加侧向作用力。
\textbf{10.} (a)--(iii)，(b)--(i)，(c)--(ii)。
\textbf{11.} \omterm{def:g10:universal-gravitation:weight}{重力} $P = 80 \times 9.81 \approx \qty{785}{N}$，阻力
为零：没有任何力与\omterm{def:g10:universal-gravitation:weight}{重力}抵消，因此速度必定改变——
跳伞者竖直向下加速。
\textbf{12.} 在阻力 $<$ \omterm{def:g10:universal-gravitation:weight}{重力}期间，作用力仍不能\omterm{def:g10:inertia:compensate}{相互抵消}：
速率持续增大，阻力也随之增大。
\textbf{13.} 匀速时：\omterm{def:g10:universal-gravitation:weight}{重力}与阻力\omterm{def:g10:inertia:compensate}{相互抵消}，因此阻力为
\qty{785}{N}。运动为\omterm{def:g10:inertia:uniform}{匀速直线运动}——一个下落物体处于
与人行道上行李箱完全相同的力学状态：
伪装的惯性。
\textbf{14.} 此时阻力大于\omterm{def:g10:universal-gravitation:weight}{重力}：作用力不能\omterm{def:g10:inertia:compensate}{相互抵消}，因此
速度改变——跳伞者仍在下降，但减速；
随着速率减小，阻力逐渐回落至\omterm{def:g10:universal-gravitation:weight}{重力}大小。
\textbf{15.} 再次匀速：阻力 $= P = \qty{785}{N}$——与
速率为 \qty{50}{m/s} 时完全相同。
\textbf{16.} 终端阻力始终等于\omterm{def:g10:universal-gravitation:weight}{重力}，无论是否张开伞衣。
伞改变的是阻力达到 \qty{785}{N} 时的\emph{速率}：
为 \qty{5}{m/s}，而非 \qty{50}{m/s}。
\textbf{17.} “\dots 作用在物体上的力\omterm{def:g10:inertia:compensate}{相互抵消}
（或无力作用）”——\cref{thm:g10:inertia:principle}的
两种读法。静止即 $\vect v = \vect 0$，只是一种普通的\omterm{def:g10:inertia:uniform}{匀速直线运动}。
\textbf{18.} 冰场：冰球滑行时身后看不到任何推动者。伽利略
指向越来越光滑的表面：使物体停下来的，是\omterm{def:g10:inertia:inventory}{摩擦力}，
而不是自然本身。
\textbf{19.} 力确实存在——\omterm{def:g10:universal-gravitation:weight}{重力} \qty{785}{N} 向下，阻力
\qty{785}{N} 向上；二者\emph{\omterm{def:g10:inertia:compensate}{相互抵消}}，这正是
速率恒定的原因。
\textbf{20.} 人行道：\omterm{def:g10:universal-gravitation:weight}{重力}与\omterm{def:g10:inertia:inventory}{法向力}。冰场：\omterm{def:g10:universal-gravitation:weight}{重力}与
\omterm{def:g10:inertia:inventory}{法向力}。跳伞：\omterm{def:g10:universal-gravitation:weight}{重力}与阻力。唯一的规律：
惯性原理，及其正、逆两种读法。
\end{solution}
